\sgaStableTarget{sga6.e05.cor.4-15}\textbf{Corollary 4.15.} \uline{Let} @@M0@@ \uline{be the} @@M1@@-\uline{ring generated over} @@M2@@ \uline{by the indeterminate} @@M3@@ \uline{subject to the relations} @@M4@@ \uline{for} @@M5@@, \uline{i}.\uline{e}. @@M6@@ (4.9). \uline{Then the} @@M7@@-\uline{ring filtration of} @@M8@@ \uline{is identical to the weight filtration with respect to the} @@M9@@.

Taking account of Proposition 4.12, Remark 4.13.ii), and the fact that @@M10@@, we have @@M11@@, so @@M12@@ is the ideal generated by the monomials of weight @@M13@@ with respect to the @@M14@@.

\sgaStableTarget{sga6.e05.sec.5}\section*{5. The @@M15@@.}

This section will be used only in Exp.VII~4, for the study of regular immersions; it may therefore be omitted on first reading.

\sgaStableTarget{sga6.e05.para.5-1}5.1. Let @@M16@@ be an integer @@M17@@. We denote by @@M18@@ the @@M19@@-@@M20@@-algebra generated by the generators @@M21@@, @@M22@@ subject to the relations @@M23@@ for @@M24@@; we therefore have:
@@M25@@

We put \sgaStableTarget{sga6.e05.eq.5-1-1}@@M26@@.

\sgaStableTarget{sga6.e05.lem.5-2}\textbf{Lemma 5.2.} \uline{For every} @@M27@@, @@M28@@ \uline{is divisible by} @@M29@@.

We know (2.3) that @@M30@@ is an isobaric polynomial of weight @@M31@@ with respect to the @@M32@@; it therefore suffices to prove the lemma for the @@M33@@.

Let @@M34@@ be the @@M35@@-ring generated over @@M36@@ by the indeterminates @@M37@@, subject to the relations @@M38@@ if @@M39@@. It is clear that @@M40@@ for @@M41@@. We can therefore define a @@M42@@-homomorphism @@M43@@ from @@M44@@ to @@M45@@ by sending @@M46@@ to @@M47@@; we then have @@M48@@,
